% Part: lambda-calculus
% Chapter: church-rosser
% Section: parallel-beta-eta-reduction

\documentclass[../../../include/open-logic-section]{subfiles}

\begin{document}

\olfileid{lam}{cr}{pbe}

\olsection{સમાંતર $\beta\eta$-લઘુકરણ}

આ વિભાગમાં $\beta\eta$-લઘુકરણને અનુરૂપ સમાંતર
$\beta\eta$-લઘુકરણ માટે ચર્ચ--રોસર ગુણધર્મ
સાબિત કરવાનો હેતુ છે.

\begin{defn}[સમાંતર $\beta\eta$-લઘુકરણ, $\beredpar$] \ollabel{defn:beredpar}
  પદો પરનું \emph{સમાંતર $\beta\eta$-લઘુકરણ}
  ($\beredpar$) આગમનથી આ રીતે વ્યાખ્યાયિત થાય છે:
  \begin{enumerate}
    \item \ollabel{defn:beredpar1} $x \beredpar x$.
    \item \ollabel{defn:beredpar2} જો $N \beredpar N'$,
      તો $\lambd[x][N] \beredpar \lambd[x][N']$.
      \sourcecorrection{OLLAM-039}{સ્થિર અંગ્રેજી મૂળ
      આ અમૂર્તીકરણ-નિયમમાં $N \xrightarrow{\beta} N'$
      લખે છે; પછીની સ્વવાચકતા અને આગમનાત્મક
      સાબિતીઓ માટે $N \beredpar N'$ જરૂરી છે.
      અગાઉના સમાંતર $\beta$-નિયમની જેમ અહીં સંબંધ
      સુધાર્યો છે.}
    \item \ollabel{defn:beredpar3} જો $P \beredpar P'$ અને
      $Q \beredpar Q'$, તો $PQ \beredpar P'Q'$.
    \item \ollabel{defn:beredpar4} જો $N \beredpar N'$ અને
      $Q \beredpar Q'$, તો
      $(\lambd[x][N])Q \beredpar \Subst{N'}{Q'}{x}$.
    \item \ollabel{defn:beredpar5} જો $N \beredpar N'$,
      તો $\lambd[x][Nx] \beredpar N'$, જો
      $x \notin FV(N)$.
  \end{enumerate}
\end{defn}

\begin{thm}\ollabel{thm:refl}
  $M \beredpar M$.
\end{thm}

\begin{proof}
  અભ્યાસપ્રશ્ન.
\end{proof}

\begin{prob}
  \olref[lam][cr][pbe]{thm:refl} સાબિત કરો.
\end{prob}

\begin{defn}[$\beta\eta$-પૂર્ણ વિકાસ]\ollabel{defn:becd}
  $M$નો \emph{$\beta\eta$-પૂર્ણ વિકાસ}~$\becd{M}$
  નીચે મુજબ વ્યાખ્યાયિત થાય છે:
  \begin{align}
    \becd{x} &= x \ollabel{defn:becd1} \\
    \becd{(\lambd[x][N])} &= \lambd[x][\becd{N}] \ollabel{defn:becd2}\\
    \becd{(PQ)} &= \becd{P}\becd{Q} && \text{જો $P$ એ $\lambd$-અમૂર્તીકરણ ન હોય}
    \ollabel{defn:becd3} \\
    \becd{((\lambd[x][N])Q)} &= \Subst{\becd{N}}{\becd{Q}}{x}
                             \ollabel{defn:becd4} \\
    \becd{(\lambd[x][Nx])} &= \becd{N} \ollabel{defn:becd5} & \text{જો $x
                                                      \notin FV(N)$}
  \end{align}
  પાંચમી કલમ લાગુ પડતી હોય ત્યારે તેને બીજી કલમથી
  પહેલાં લાગુ કરવાની છે; બીજી કલમ બાકી અમૂર્તીકરણો
  માટે છે. નહિ તો એક જ $\eta$-સંકોચ્ય પદને બંને
  કલમો જુદાં મૂલ્યો આપે.
  \sourcecorrection{OLLAM-040}{સ્થિર અંગ્રેજી મૂળની
  બીજી કલમ દરેક અમૂર્તીકરણ માટે છે, જ્યારે પાંચમી
  કલમ કેટલાક અમૂર્તીકરણો માટે જુદું પરિણામ આપે છે.
  આથી વ્યાખ્યા પરસ્પર આવરી લેતી અને અસંગત બને છે.
  અહીં પાંચમી કલમને પ્રાથમિકતા આપી બીજી કલમને
  બાકી કિસ્સાઓ સુધી મર્યાદિત કરી છે; આ પસંદગીની
  અનુરૂપ સાબિતી આગળ હજી જરૂરી છે.}
\end{defn}

\begin{lem}\ollabel{lem:comp}
  $M \beredpar M'$ અને $R \beredpar R'$ હોય તો
  $\Subst{M}{R}{y} \beredpar
  \Subst{M'}{R'}{y}$.
\end{lem}

\begin{proof}
  $M \beredpar M'$ની !!{derivation} પર આગમનથી.

  પહેલા ચાર કિસ્સા \olref[pb]{lem:comp} જેવા છે;
  ત્યાં નોંધાયેલો સ્થાનાપન્ન-સંયોજનનો સાબિતી-ગાળો
  અહીં પણ રહે છે. છેલ્લો નિયમ
  \olref{defn:beredpar5} હોય તો
  $M$ એ $\lambd[x][Nx]$ અને $M'$ એ $N'$, કોઈક
  $x$ અને~$N'$ માટે, જ્યાં $x \notin FV(N)$ અને
  $N \beredpar N'$. બતાવવું છે કે
  $\Subst{(\lambd[x][Nx])}{R}{y} \beredpar
  \Subst{N'}{R'}{y}$. $x$ને $y$, $R$ અને~$R'$થી
  તાજો પસંદ કરીએ તો જરૂરી દાવો
  $\lambd[x][\Subst{N}{R}{y} x] \beredpar
  \Subst{N'}{R'}{y}$ બને છે. અહીં
  $x \notin FV(\Subst{N}{R}{y})$ પણ છે.
  તેથી આગમનની ધારણા અને
  \olref{defn:beredpar}\olref{defn:beredpar5}
  લાગુ પડે છે.
  \sourcecorrection{OLLAM-041}{સ્થિર અંગ્રેજી મૂળ
  $\eta$-કિસ્સામાં સ્થાનાપન્ન પછીનો બાંધનાર~$x$
  તાજો રહે છે એમ કોઈ શરત વિના માને છે. વર્ગના
  પ્રતિનિધિમાં~$x$ને $y$, $R$ અને~$R'$થી તાજો
  પસંદ કરીને જ પાંચમો નિયમ લગાડી શકાય. પહેલા
  ચાર કિસ્સા પણ \olref[pb]{lem:comp}ના અગાઉના
  ખુલ્લા સંયોજન-ગાળાને વારસામાં લે છે.}
\end{proof}

\begin{lem}\ollabel{lem:cont}
  $M \beredpar M'$ હોય તો
  $M' \beredpar \becd{M}$.
\end{lem}

\begin{proof}
  $M \beredpar M'$ની !!{derivation} પર આગમનથી.

  સ્થિર મૂળ પહેલા ચાર કિસ્સા
  \olref[pb]{lem:cont} જેવા કહે છે. પરંતુ
  $M$ એ $\eta$-સંકોચ્ય અમૂર્તીકરણ હોય ત્યારે
  ઉપર પ્રાથમિકતા આપેલી પાંચમી કલમથી
  $\becd{M}$ નક્કી થાય છે; સામાન્ય
  અમૂર્તીકરણ-કિસ્સાની જૂની દલીલ એ પરિણામ સુધી
  પહોંચે છે એવું અહીં સાબિત થયું નથી.
  \sourcecorrection{OLLAM-042}{સ્થિર અંગ્રેજી મૂળની
  પૂર્ણ-વિકાસ વ્યાખ્યાના બીજા અને પાંચમા નિયમો
  અથડાય છે. પાંચમાને પ્રાથમિકતા આપ્યા પછી
  અમૂર્તીકરણ હેઠળના સમાંતર પગલાં માટે સંપૂર્ણ
  વિકાસ સુધીનું એક વધુ પગલું મળે છે તે અલગ
  સાબિત કરવું પડે; ``પહેલા ચાર કિસ્સા અગાઉ જેવા''
  કહેવું પૂરતું નથી. આ સાબિતી-ગાળો ખુલ્લો છે.}
  છેલ્લો નિયમ \olref{defn:beredpar5} હોય તો
  $M$ એ $\lambd[x][Nx]$ અને $M'$ એ $N'$, કોઈક
  $x$, $N$, $N'$ માટે, જ્યાં $x \notin FV(N)$ અને
  $N \beredpar N'$. બતાવવું છે કે
  $N' \beredpar \becd{(\lambd[x][Nx])}$,
  એટલે પાંચમી કલમ પ્રમાણે
  $N' \beredpar \becd{N}$. આ આગમનની ધારણાથી
  તરત મળે છે, પરંતુ ઉપરનો અમૂર્તીકરણ-કિસ્સો
  હજી ખુલ્લો રહે છે.
\end{proof}

\begin{thm}\ollabel{thm:cr}
  $\beredpar$ ચર્ચ--રોસર ગુણધર્મ ધરાવે છે.
\end{thm}

\begin{proof}
  \olref{lem:cont}માંથી આ પરિણામ મળે, જો તેની
  ઉપર નોંધેલી સાબિતી-ખામી પૂરી કરવામાં આવે.
\end{proof}

\end{document}
